{\renewcommand{\textfraction}{0.01}
\renewcommand{\bottomfraction}{0.99}
\renewcommand{\floatpagefraction}{0.99}
\chapter{Rubrics, Prompts, and Full Task Examples}
\label{app:rubrics-prompts}

\vspace{-2.5em}
\section{Reconstructed One-Shot PF Generation Prompt}
\label{app:oneshot-prompt}
\vspace{-0.4em}
Table~\ref{tab:app-oneshot-prompt} presents the reconstructed six-part prompt used in Chapter~\ref{chap:oneshot}. The system prompt stays fixed across runs, while the user message supplies the target concept, theme, interest, and prerequisites for each grade~7 cell.

\begin{table}[!htbp]
\centering
\vspace{-0.5em}
\caption[Reconstructed six-part one-shot prompt for PF task generation]{Reconstructed six-part one-shot prompt for PF task generation.}
\label{tab:app-oneshot-prompt}
\vspace{0.1em}
{\fontsize{7.5}{8.6}\selectfont
\setlength{\tabcolsep}{4pt}
\renewcommand{\arraystretch}{1.03}
\begin{tabular}{@{} >{\raggedright\arraybackslash}p{0.15\textwidth} >{\raggedright\arraybackslash}p{0.825\textwidth} @{}}
\toprule
\textbf{Component} & \textbf{Prompt Content} \\
\midrule
\textbf{1. Role \& Rules}\newline (System) &
\texttt{ROLE AND OBJECTIVE:} Educational design assistant creating a mathematics PF problem for the supplied grade, target, theme, student interest, and explicitly listed prior knowledge.\newline
\texttt{RESPONSE RULES:} Write an engaging, age-appropriate narrative with the needed data and an open-ended challenge. Leave room for students to choose and explain approaches. Do not supply a worked solution, teacher commentary, or steps teaching the target concept. Output only JSON: \texttt{\{"status": "ready", "title": "...", "body": "...", "question": "..."\}}. \\
\midrule
\textbf{2. PF Overview}\newline (System) &
\texttt{PRODUCTIVE FAILURE OVERVIEW:} Students attempt a challenging problem before direct instruction on the target concept so their initial representations and approaches support later comparison and consolidation. Support familiar knowledge while making further conceptual development useful. \\
\midrule
\textbf{3. PF Rubric}\newline (System, \cite{rahimi2025}) &
\texttt{PUBLISHED RUBRIC:} Use these five dimensions (1--3 anchors) as design guidance without self-scoring:\newline
\textbf{1.~APK:} 1: No prerequisite concepts; 2: One prerequisite; 3: Two or more prerequisites.\newline
\textbf{2.~SSC:} 1: Too easy/difficult (disengagement); 2: Moderate challenge with some failures; 3: Challenging but manageable, leading to productive engagement.\newline
\textbf{3.~AE:} 1: Not relatable/meaningful; 2: Somewhat engaging, realistic; 3: Highly engaging narrative.\newline
\textbf{4.~MR:} 1: Few/no methods; 2: Limited diversity; 3: Wide range of representations and methods.\newline
\textbf{5.~OE:} 1: Simple, well-structured, single path; 2: Moderate complexity with limited variables; 3: Highly complex, integrating multiple variables and representations. \\
\midrule
\textbf{4. Five Literature Examples}\newline (System, \textsc{StratL}~\cite{puech2025}) &
\texttt{FIVE LITERATURE EXAMPLES:} Concise adaptations of human-source tasks from Puech et al.\ (2025) illustrating PF task forms; do not copy their topic or data when unmatched to the target.\newline
\textbf{Ex.~1 -- Consistency (Kapur, 2015):} Five-season goal totals for Mike (13, 12, 15, 17, 13), Dave (12, 14, 16, 15, 13), and Ivan (14, 10, 18, 18, 15); design consistency measures using all five seasons.\newline
\textbf{Ex.~2 -- Representation (Kerrigan, 2018):} Six states have populations of 1,646k; 6,936k; 154k; 2,091k; 685k; and 988k, with 250 legislative seats; propose and justify a fair seat allocation.\newline
\textbf{Ex.~3 -- Land sharing (Kerrigan et al., 2021):} Four owners value four plots at Andrea (150, 300, 200, 350), Beatrice (200, 350, 100, 350), Cleo (200, 450, 200, 150), David (250 each); justify a fair division.\newline
\textbf{Ex.~4 -- Shared furniture (Kerrigan et al., 2021):} Three roommates value four items at Ashley (150, 150, 175, 275), Bill (250, 200, 150, 300), Cliff (175, 125, 125, 250); justify a fair settlement.\newline
\textbf{Ex.~5 -- Sharing a bicycle (Kapur, 2010):} Two friends have 5~km left and one bike (Jasmine: 150~m/min ride, 30~m/min walk; Hady: 200~m/min ride, 50~m/min walk); find where to leave it to tie. \\
\midrule
\textbf{5--6. Topic Cell Specification}\newline (User, Grade~7) &
Generate one paper-setting PF task as JSON for the supplied cell (target concept is untaught):\newline
\textbf{Cell 1 -- Sports (Basketball):} \emph{Target:} Recognize/represent proportional relationships. \emph{Interest:} Basketball and Stephen Curry. \emph{Prerequisites:} Fractional unit rates; rational addition/subtraction.\newline
\textbf{Cell 2 -- Shopping (Pocket Money):} \emph{Target:} Multistep ratio/percent problems (tax, markups, markdowns). \emph{Interest:} Shopping and pocket money. \emph{Prerequisites:} Proportions; rational operations.\newline
\textbf{Cell 3 -- Music (Playlists):} \emph{Target:} Compare two populations via center/variability from random samples. \emph{Interest:} Bands, songs, playlists. \emph{Prerequisites:} Mean, median, IQR, MAD; decimals.\newline
\textbf{Cell 4 -- Home \& Garden:} \emph{Target:} Area of composite 2D figures (triangles, quadrilaterals, polygons). \emph{Interest:} Garden, patio, or mural design. \emph{Prerequisites:} Composing/decomposing area; rationals. \\
\bottomrule
\end{tabular}
}
\end{table}
\clearpage

\section{Automated Five-Dimension Rubric Judge Prompt}
\label{app:oneshot-judge}
\vspace{-0.4em}
Table~\ref{tab:app-oneshot-judge} presents the prompt used for the two-pass blinded LLM rubric evaluation in Chapter~\ref{chap:oneshot}. All five dimension names and 1--3 scoring anchors are retained verbatim from Table~1 of Rahimi et al.~\cite{rahimi2025}, while the input specification supplies the shared topic-cell setting and student-facing task text without revealing the generating model.

\begin{table}[!htbp]
\centering
\vspace{-0.4em}
\caption[Automated five-dimension PF rubric judge prompt]{Automated five-dimension PF rubric judge prompt adapted from Rahimi et al.~\cite{rahimi2025}.}
\label{tab:app-oneshot-judge}
\vspace{0.1em}
{\fontsize{7.5}{8.8}\selectfont
\setlength{\tabcolsep}{4pt}
\renewcommand{\arraystretch}{1.06}
\begin{tabular}{@{} >{\raggedright\arraybackslash}p{0.17\textwidth} >{\raggedright\arraybackslash}p{0.805\textwidth} @{}}
\toprule
\textbf{Component} & \textbf{Prompt Content} \\
\midrule
\textbf{1. Evaluator Role}\newline (System) &
Evaluate a mathematics Productive Failure problem using the five dimensions in Rahimi et al.\ (2025), Table~1. The scoring anchors below are the paper's verbatim anchors. This is an automated rating, not a human rating or observed student learning outcome. \\
\midrule
\textbf{2. Verbatim Rubric Anchors}\newline (System, \cite{rahimi2025}) &
\textbf{1.~Activation of Prior Knowledge (APK):} 1~(Low): No prerequisite concepts; 2~(Moderate): Includes one prerequisite concept; 3~(High): Includes two or more prerequisite concepts.\newline
\textbf{2.~Sweet Spot Calibration (SSC):} 1~(Low): Too easy or too difficult: disengagement; 2~(Moderate): Moderate challenge with some failures; 3~(High): Well-calibrated: challenging but manageable, leading to productive engagement.\newline
\textbf{3.~Affective Engagement (AE):} 1~(Low): Not relatable or meaningful: fails to engage; 2~(Moderate): Somewhat engaging with relatable/realistic elements; 3~(High): Highly engaging, presented in meaningful narrative or real-world scenario.\newline
\textbf{4.~Multiple Representations \& Solutions (MR):} 1~(Low): Few or no representations or solution methods; 2~(Moderate): Some diversity but limited; 3~(High): Wide range of representations and solution methods; fosters creativity \& collaboration.\newline
\textbf{5.~Open-endedness / Ill-structuredness (OE):} 1~(Low): Simple, well-structured, single path; 2~(Moderate): Moderate complexity with limited variables; 3~(High): Highly complex, requiring integration of multiple variables and representations. \\
\midrule
\textbf{3. Input Scope \& Reference Rules}\newline (System) &
\texttt{INPUT:} \texttt{task\_setting} states the intended PF exploration use, the untaught learning target, and the assigned theme and student interest. \texttt{learner\_context} lists the assumed prior knowledge. \texttt{student\_task} is the exact text students receive and is the only item rated. Use \texttt{task\_setting} and \texttt{learner\_context} only as the reference a trained rater would know: prerequisites for APK are the listed prior knowledge, challenge for SSC is for students who know that prior knowledge but not the target, and relatability for AE is for students with the stated interest. Do not reward a problem for restating the setting, and do not rate the setting itself. \\
\midrule
\textbf{4. Rating \& Output Contract}\newline (System) &
Apply each dimension independently and give brief, concrete evidence from the supplied problem and learner context before each score. Do not target a score distribution, match published model averages, or add extra scoring rules (such as requiring named characters, deliberately wrong answers, fixed counts of display formats, or banning a unique numerical answer). Assess what the problem supports rather than merely what it claims. If evidence is insufficient, return \texttt{null} for its score and explain why; keep mathematical concerns in a separate list without inventing a sixth score. Return only JSON with keys \texttt{APK}, \texttt{SSC}, \texttt{AE}, \texttt{MR}, \texttt{OE} (each \texttt{\{"evidence": "...", "score": 1|2|3|null\}}), and \texttt{mathematical\_concerns} (\texttt{[...]}). Code computes the 5--15 total. \\
\midrule
\textbf{5. Blinded Task Payload}\newline (User) &
JSON object containing \texttt{task\_setting} (\texttt{intended\_use}, \texttt{learning\_target}, \texttt{theme}, \texttt{student\_interest}), \texttt{learner\_context} (\texttt{grade}, \texttt{subject}, \texttt{jurisdiction}, \texttt{assumed\_prior\_knowledge}), and \texttt{student\_task} (concatenated student-facing \texttt{title}, \texttt{body}, and \texttt{question}), shuffled under a blind task ID with all model and condition labels removed. \\
\bottomrule
\end{tabular}
}
\end{table}
\clearpage
}

